\section{Introduction}\label{sec:intro}

Modern Linux filesystems detect metadata corruption through cryptographic
or polynomial checksums; few correct it in place, and to the
author's knowledge none ship with a \emph{formally proved} discipline
for recovery from \emph{active} memory corruption induced by physical
perturbation.\footnote{FSCQ~\cite{chen2015fscq} and
Yggdrasil~\cite{sigurbjarnarson2016yggdrasil} prove crash-consistency,
that is, recovery under power-loss with a benevolent storage substrate.
Neither models adversarial corruption of volatile memory under physical
perturbation, which is the subject of the present work.}
Detection without proved correction is insufficient for radiation-prone
deployments \cite{nasa-seeca,esa-eccs60,normand1996}: a single-event
upset (SEU) on a critical metadata word can put the filesystem into a
state from which best-effort recovery procedures diverge, loop, or
silently propagate the corruption.

The companion Technical Report~\cite{desbrieres2026ftrfs} documents an
implementation, FTRFS, that adds layered Reed--Solomon forward error
correction to inode, allocation-bitmap, and superblock metadata, with a
persistent radiation event journal. FTRFS as a name and concept
originates in the work of Fuchs \emph{et al.}~\cite{fuchs2015ftrfs};
the companion report builds on that lineage with contemporary Linux
kernel infrastructure. That implementation reaches a
ceiling: it demonstrates correction \emph{empirically} but does not
characterize, prove, or bound its recovery properties. The present
report addresses that gap by treating recovery itself as a mathematical
object, and by tracing it explicitly to a kernel implementation path.

\paragraph{Position of the problem.}
Let~\(\mathcal{S}\) denote the joint state space of a filesystem's
volatile (RAM, page cache, kernel buffers) and persistent (NVMe,
NAND-flash) memory. Let~\(\mathcal{E}\) denote the family of physically
realizable Single-Event Effects acting on~\(\mathcal{S}\), parameterized
by the linear energy transfer of the incident particle, the vulnerable
cross-section of the target memory cell, and the geometry of the
silicon node. Let~\(\models\) denote the relation ``\(s \in
\mathcal{S}\) is consistent'' (cyclic redundancy code valid, Reed--Solomon
syndromes zero, structural sentinels intact, on-disk versus in-memory
copies coherent). The problem is to construct an
operator~\(\mathcal{G} : \mathcal{S} \to \mathcal{S} \cup \{\bot\}\),
implementable in a Linux-kernel filesystem, such that
\[
  \forall\, s_0 \models,\ \forall\, e \in \mathcal{E},\
  \exists\, k \in \mathbb{N} :\
  \mathcal{G}^{k}\bigl(e(s_0)\bigr) \models
  \quad\lor\quad
  \mathcal{G}^{k}\bigl(e(s_0)\bigr) = \bot,
\]
where~\(\bot\) denotes the fail-closed state in which the filesystem
explicitly refuses to proceed. The disjunction is essential: silent
corruption is excluded by construction.

\paragraph{Where we are going, and how.}
This report defends a single thesis: \emph{provable filesystem recovery
under SEE is achievable, and the chain from theorem to kernel code is
auditable line by line}. The chain is presented in
\cref{sec:m2c} as a three-tier descent---mathematics, algorithm,
kernel implementation---closed by a traceability table
(\cref{tab:traceability}) where each line of the proof is paired with
the algorithmic step and the kernel construct that realizes it. This
structure is the report's central contribution to the literature: it
makes the implementation gap explicit and bounded.

\paragraph{Contributions.}
We claim the following:
\begin{itemize}
\item A multi-layer physical threat model
  (\cref{sec:phys}) drawing on canonical SEE literature
  \cite{heinrich1977,petersen1996,dodd2003,baumann2005,normand1996} and
  standards~\cite{jedec89a,nasa-hdbk-4002a,esa-eccs60}, treating DRAM,
  SRAM, and NAND-flash distinctly with realistic per-technology rates.
\item A formal definition of the filesystem state space, its
  consistency relation, and a small set of structural invariants
  (\cref{sec:state}) that are necessary preconditions for any proved
  recovery.
\item A soundness theorem (\cref{thm:soundness}, full proof in
  \cref{app:proof}) showing that under stated hypotheses on the
  perturbation family, the recovery operator drives every initial
  consistent state to a consistent successor or to fail-closed in a
  finite number of iterations.
\item A three-tier descent (\cref{sec:m2c}) from the operator's
  mathematical definition to a deterministic algorithm to a Linux-kernel
  implementation path, with a traceability table
  (\cref{tab:traceability}) auditable line by line. This is the
  report's central contribution.
\item An empirical re-reading (\cref{sec:empirical}) of the FTRFS v1
  bench dataset~\cite{desbrieres2026ftrfs} through the recovery-calculus
  lens, identifying which metrics constrain recovery and which do not.
\end{itemize}

\paragraph{Scope.}
This is v1. We model SEE explicitly; we do not model Total Ionizing
Dose (TID) or Displacement Damage Dose (DDD), both of which are
hardware-level cumulative phenomena outside the filesystem's
operational reach. We discuss DDR5 on-die ECC~\cite{jedec79-5} as a
recent evolution that modifies the threat surface, but we do not
incorporate it into the v1 model: the matrix-cell-level abstraction
remains the conservative baseline, and on-die ECC will be added in v2
as a strengthening assumption (it can only improve outcomes, never
worsen them, given the soundness theorem's structure). We claim no
beam-time measurements; the empirical section reuses the FTRFS v1
software validation dataset and projects expected outcomes from
canonical cross-section data. The recovery calculus is positioned
to support the auditable-evidence posture demanded by safety-critical
standards such as MIL-STD-882E~\cite{milstd882e} and DO-178C, which
require traceable failure-mode coverage; the three-tier descent of
\cref{sec:m2c} is designed to provide such a trace.

\paragraph{Anti-claim.}
We are explicit about what BEAMFS v1 does \emph{not} prove:
filesystem behavior under simultaneous TID-induced cell drift and
SEE; recovery from latched-up hardware states (SEL); coverage of
adversarial fault injection by an attacker with kernel privileges
(distinct from physical SEE). Each is acknowledged as future work
(\cref{sec:limits}).
